\documentclass[letterpaper,11pt]{article}
\usepackage{graphicx}
\usepackage{geometry}
\linespread{1.2}                        
\setlength{\oddsidemargin}{0.0in}            
\setlength{\textwidth}{6.5in}    
\usepackage[utf8]{inputenc} % Asegúrate de usar utf8
\usepackage[T1]{fontenc} % Asegúrate de usar T1 font encoding      
\setlength{\topmargin}{-.6in}            
\setlength{\textheight}{9in}
\usepackage{latexsym} 
\usepackage{amssymb} 
\usepackage{amsmath}
\usepackage{amsxtra}
\usepackage[mathcal]{eucal} 
\usepackage{enumerate}
\usepackage{hyperref}
\usepackage{rotating}
\usepackage{caption}
\usepackage{lscape}
\usepackage{paralist} 
\usepackage{indentfirst}
\usepackage[dvipsnames,svgnames]{xcolor} 
\usepackage{setspace}
\usepackage{multicol} 
\usepackage{float}

\definecolor{codegreen}{rgb}{0,0.6,0}
\definecolor{codegray}{rgb}{0.5,0.5,0.5}
\definecolor{codepurple}{cmyk}{0.65, 1, 0, 0.35}
\definecolor{backcolour}{rgb}{0.95,0.95,0.92}

\usepackage{listings}
\usepackage{xcolor}

\lstset{
    language=R, % lenguaje de programación
    basicstyle=\small\ttfamily, % estilo básico
    commentstyle=\color{gray}, % estilo para los comentarios
    keywordstyle=\color{blue}, % estilo para las palabras clave
    stringstyle=\color{black}, % estilo para las cadenas de texto
    showstringspaces=false, % no mostrar espacios en cadenas de texto
    tabsize=4, % tamaño de la tabulación
    breaklines=true, % permitir el salto de línea automático
    postbreak=\mbox{\textcolor{red}{$\hookrightarrow$}\space}, % símbolo al final de una línea partida
    numbers=left, % mostrar números de línea
    numberstyle=\tiny\color{gray}, % estilo para los números de línea
    frame=single, % mostrar un borde alrededor del código
    backgroundcolor=\color{gray!10}, % color de fondo
}

\begin{document}

\begin{center}
    {\Large  Taller 2: Econometría }

Profesora:    Jocelyn Tapia Stefanoni\\
Ayudantes:   Matías Balboa y   Fabián Cisternas
\end{center}

\textbf{Instrucciones:} 

\begin{itemize}
    \item Usted dispone de 65 minutos  para intentar los 100 puntos del taller.
    \item Cuide la presentación y redacción de sus respuestas.
    \item Puede utilizar su computador y los apuntes tanto de clase como de ayudantía.
\item Descargue sus respuestas en un archivo con la extensión .ipynb. Debe enviar el enlace al notebook en Google Colab junto con el archivo .ipynb por AULA. 
    \end{itemize}


Considere los datos del archivo \texttt{elem94\_95} del paquete \texttt{wooldridge} y responda las siguientes preguntas. 
 



\begin{enumerate}
\item (8 puntos) Describa la base datos, señalando cuál es su origen y que variables presenta.

\item (8 puntos) Utilizando \texttt{stargazer} realice una tabla con la estadística descriptiva de las variables:  \texttt{lavgsal}
\texttt{bs},\texttt{lenrol}, \texttt{lstaff} y \texttt{lunch}.  

    \item  (18 puntos) Realice una regresión lineal simple de \texttt{lavgsal}, que representa el logaritmo del sueldo promedio del profesor, sobre \texttt{bs} que corresponde al cociente de beneficios promedio entre sueldo promedio (por escuela). 
    Es la pendiente estimada estadísticamente distinta de cero? 
    Es estadísticamente diferente de -1?

    \item (14 puntos) Construya un intervalo de confianza para el parámetro de la pendiente? Explique como se relacionan  sus resultados con los del item anterior.

    \item (12 puntos) Agregue las variables \texttt{lenrol} (logaritmo de la matrícula escolar) y \texttt{lstaff} (logaritmo de la cantidad de personal) 
    a la regresión del inciso anterior. 
    Qué ocurre con el coeficiente de \texttt{bs}? Para realizar su comparaci\'on muestre los resultados de ambas regresiones en una única  tabla utilizando \texttt{stargazer}.
   
    % \item ¿A qué se debe que el error estándar del coeficiente de \texttt{bs} sea menor en el inciso (ii) que en el inciso (i)? 
    % \textit{Sugerencia:} analice qué ocurre con la varianza del error frente a la multicolinealidad 
    % cuando se agregan \texttt{lenrol} y \texttt{lstaff}.

    \item (18 puntos) Testee que la elasticidad  del salario respecto de la cantidad de personal sea  10 veces la elasticidad del salario respecto de la  matr\'icula escolar.
    

    \item (10 puntos) Ahora agregue a la regresión la variable \texttt{lunch} (porcentaje de alumnos que reciben desayuno gratuito). Controlando por los demás factores, tiene un efecto estadísticamente significativo  enseñar a estudiantes con este beneficio?
   \item (12 puntos) Compare los resultados de sus 3 estimaciones utilizando \texttt{stargazer}.

\end{enumerate}
\end{document}